\section{Mô phỏng thế giới}

\subsection{Câu hỏi nghiên cứu}

\textit{Character có thể tồn tại trong một world persistent thay vì chỉ phản hồi từng message?}

\subsection{Evidence}

\begin{table}[H]
\centering
\caption{So sánh các Hệ thống World Simulation}
\label{tab:world-sim}
\begin{tabular}{lllll}
\toprule
\textbf{System} & \textbf{Agents} & \textbf{Fidelity} & \textbf{Scalability} & \textbf{Thành tựu chính} \\
\midrule
Voyager & 1 & Cao (3.3$\times$ items, 56 skills) & Thấp & Embodied learning \\
CharacterBox & Variable & Trung bình (78\% consistency @200 turns) & Trung bình & Persistent persona \\
GenSim & 100,000 & Thấp (depth 3/10) & Xuất sắc & Scale record \\
CAREB-MAS & Variable & Cao & Trung bình & 5 emergent social phenomena \\
\bottomrule
\end{tabular}
\end{table}

\subsection{Tradeoff giữa Scalability và Fidelity}

Hệ scale được (GenSim 100K) thì character nông (depth 3/10). Hệ sâu (Voyager) thì chỉ 1 agent. \textbf{Chưa có solution cân bằng.}

\subsection{Khuyến nghị}

\textbf{Hybrid approach}: Structured core (JSON state) + narrative overlay (text generation). Aivora Lab nên theo hướng này với:
\begin{itemize}
    \item External memory store (SQLite)
    \item Character state serialization
    \item Spatial partitioning cho multi-agent
\end{itemize}

% =============================================================================
% 12. MULTI-AGENT